%% kapa: exam=1 round=20 year=2009 period=1 form=A subject=civil_law subject_ko=민법 seq=10 number=10
%% source: Q-Net 감정평가사 기출문제 (cst00309) artlSeq=5024178 file=extracted/1008/5024178/2009년 제20회 감정평가사 1차시험 기출문제/1교시.hwp
%% conversion: kordoc 4.24.51 -> tools/gen1.py (auto); review=proofread-sonnet; flags=none
\begin{question}{10}
물권의 취득을 위하여 등기를 요하지 \uline{않는} 경우를 모두 묶은 것은?
\begin{framedtext}
ㄱ. 경매로 인한 소유권의 취득 \par
ㄴ. 상속에 의한 소유권의 취득 \par
ㄷ. 법원의 이행판결에 의한 소유권의 취득 \par
ㄹ. 공용징수에 의한 소유권의 취득 \par
ㅁ. 저당물의 경매로 인한 법정지상권의 취득 \par
ㅂ. 점유취득시효에 의한 소유권의 취득 \par
\end{framedtext}
\begin{choices}
  \item ㄱ, ㄴ
  \item ㄷ, ㅂ
  \item ㄱ, ㄴ, ㄷ
  \item ㄱ, ㄴ, ㄹ, ㅁ
  \item ㄷ, ㄹ, ㅁ, ㅂ
\end{choices}
\end{question}
\answer{4}
